% Preambule

% Nastavení tabulky
\renewcommand{\arraystretch}{0.7} 
\setlength{\extrarowheight}{15pt}  
\setlength{\tabcolsep}{10pt}
\newgeometry{nohead, nofoot, bottom=1.125cm, top=1.125cm, left =1.25cm, right=1.25cm,margin=0cm,twoside=false}


% Tabulka
\small % menší písmo (platí jen pro tabulku)
\begin{xltabular}{\textwidth}{>{\raggedleft}p{2cm}X}

{\larken -100000} & Začátek rozšiřování řeči u Homo sapiens. \\
{\larken -1400} & Nejstarší dochovaná píseň: Hymnus pro Nikkalu. \\
{\larken 1945} & Na Hirošimu a Nagasaki jsou svrženy jaderné bomby „Little Boy“ a „Fat Man“, okamžitě zahynulo kolem 100 000 lidí \\
{\larken 1957} & Waldemar Matuška začíná svou kariéru v pražském klubu Reduta a brzy se stává výraznou osobností divadla Semafor. Zpívá sólově i v duetech a objevuje se také ve filmech. \\
{\larken 1960} & Spirituál kvintet vzniká z vysokoškolského souboru a postupně se stává jednou z nejdéle působících českých skupin. I přes proměnlivý počet členů si udržel název „kvintet“. \\
{\larken 1962} & Vzniká kapela Olympic (původně Karkulka), která se brzy stává jedním ze symbolů českého rocku. Název se ustálil podle klubu, kde skupina působila. \\
{\larken 1968} & Žalman působí na scéně, nejprve ve skupině Minnesengři. V roce 1982 zakládá Žalman a spol., kteří koncertují v různých sestavách dodnes. \\
{\larken 1971} & Pacifik navazuje na trampskou tradici a vzniká z bývalých členů jiných skupin. Název vymyslel Tony Linhart, který také stál u prvních písní. \\
{\larken 1974} & Wabi Daněk vydává píseň „Rosa na kolejích“, která se stává jednou z nejznámějších trampských skladeb. \\
{\larken 1979} & Robert Křesťan přechází k Poutníkům a později zakládá Druhou trávu, čímž navazuje na bluegrassovou tradici. \\
{\larken 1982} & Vzniká skupina Klíč (původně Petrklíč), inspirovaná středověkými a renesančními motivy. Ve stejném roce se formuje i Visací zámek, který přináší jednoduchý a energický punkový repertoár. \\
{\larken 1988} & Čechomor vzniká ve Svitavách a originálním způsobem zpracovává lidové písně. \\
{\larken 1992} & Premier – poprocková kapela ze Zlína založená Danielem Hrnčiříkem a Jaroslavem Bobowskim. \\
{\larken 1994} & Vypsaná fiXa vzniká v Pardubicích. Je ovlivněna kapelami Pixies a The Cure a známá svými osobitými texty. \\
{\larken 1995} & Traband zakládají tři muzikanti v čele s Jardou Svobodou. Název je odvozen od slovní hříčky „troj-bend“ a zároveň od původního významu slova „trabant“ jako průvodce či ochránce. \\
{\larken 1996} & Ready Kirken vzniká z domácích nahrávek Michala Hrůzy. Název je převzat z pseudonymu „Reddy Kirken“. \\
{\larken 1998} & UDG vzniká v Ústí nad Labem. Název je zkratkou „Useless Demi-Gods“, vytvořenou z výrazů z textů Manic Street Preachers. \\
{\larken 2001} & Vychází singl Karla Gotta „Být stále mlád“. \\
{\larken 2002} & Xavier Baumaxa začíná sólově vystupovat po předchozích zkušenostech a inspiraci jinými písničkáři. \\
{\larken 2012} & Na albu „Výhledově! Best Of 25 let“ vychází píseň „Hirošima Nagasaki“. \\
{\larken 2019} & Umírá Karel Gott. \\
{\larken 2023} & Zveřejněna platforma Suno pro generování písní pomocí umělé inteligence. \\
{\larken 2025} & OpenAI je nařízeno zaplatit odškodnění za trénování jazykových modelů na dílech umělců bez povolení. \\

\end{xltabular}

% návrat k normálu (pokud pokračuje další text)
\normalsize
\restoregeometry
